% from ega3-1-fr.tex
\oldpage[III]{89}
Pour tout $\mathcal O_X$-Module quasi-cohérent $\mathcal F$,
l'homomorphisme canonique
$R^pg_*(f_*(\mathcal F))\to R^p(g\circ f)_*(\mathcal F)$ (0, 12.2.5.1)
est bijectif pour tout $p$.
\end{corollary}

\begin{proof}
Il suffit de remarquer que, d'après
(\hyperref[III.1.3.3-fr]{1.3.3}), pour tout ouvert affine $W$ de $Z$,
l'homomorphisme canonique
$H^p(g^{-1}(W),f_*(\mathcal F))\to
H^p(f^{-1}(g^{-1}(W)),\mathcal F)$ est bijectif~; cela prouve que
l'homomorphisme de préfaisceaux définissant l'homomorphisme canonique
$R^pg_*(f_*(\mathcal F))\to R^p(g\circ f)_*(\mathcal F)$ est bijectif
(0, 12.2.5).
\end{proof}

\subsection{Application à la cohomologie des préschémas quelconques}
\label{subsection:III.1.4-fr}

\begin{proposition}[1.4.1]
\label{III.1.4.1-fr}
Soient $X$ un schéma, $\mathfrak U=(U_\alpha)$ un recouvrement de $X$ par
des ouverts affines. Pour tout $\mathcal O_X$-Module quasi-cohérent
$\mathcal F$, les modules de cohomologie $H^\bullet(X,\mathcal F)$ et
$H^\bullet(\mathfrak U,\mathcal F)$ (sur $\Gamma(X,\mathcal O_X)$) sont
canoniquement isomorphes.
\end{proposition}

\begin{proof}
En effet, comme $X$ est un schéma, toute intersection finie $V$ d'ouverts du
recouvrement $\mathfrak U$ est affine (I, 5.5.6), donc
$H^q(V,\mathcal F)=0$ pour $q\geq1$ en vertu de
(\hyperref[III.1.3.1-fr]{1.3.1}). La proposition résulte alors du th. de
Leray (G, II, 5.4.1).
\end{proof}

\begin{remark}[1.4.2]
\label{III.1.4.2-fr}
On notera que la conclusion de (\hyperref[III.1.4.1-fr]{1.4.1}) est encore
valable lorsque les intersections finies des ensembles $U_\alpha$ sont
affines, même lorsque l'on ne suppose pas nécessairement que $X$ soit un
schéma.
\end{remark}

\begin{corollary}[1.4.3]
\label{III.1.4.3-fr}
Soient $X$ un schéma dont l'espace sous-jacent est quasi-compact,
$A=\Gamma(X,\mathcal O_X)$, $\mathbf f=(f_i)_{1\leq i\leq r}$ une suite
finie d'éléments de $A$ tels que les $X_{f_i}$ (notations de
(\hyperref[III.1.2.1-fr]{1.2.1})) soient affines. Alors (avec les notations
de (\hyperref[III.1.2.1-fr]{1.2.1})), pour tout $\mathcal O_X$-Module
quasi-cohérent $\mathcal F$, on a un isomorphisme canonique fonctoriel en
$\mathcal F$
\begin{equation}
\label{III.1.4.3.1-fr}
\tag{1.4.3.1}
H^q(U,\mathcal F)\xrightarrow{\sim}H^{q+1}((\mathbf f),M)
\qquad\text{pour }q\geq1
\end{equation}
et une suite exacte fonctorielle en $\mathcal F$
\begin{equation}
\label{III.1.4.3.2-fr}
\tag{1.4.3.2}
0\longrightarrow H^0((\mathbf f),M)\longrightarrow M
\longrightarrow H^0(U,\mathcal F)
\longrightarrow H^1((\mathbf f),M)\longrightarrow0.
\end{equation}
\end{corollary}

\begin{proof}
Cela résulte en effet aussitôt de
(\hyperref[III.1.4.1-fr]{1.4.1}) et
(\hyperref[III.1.2.3-fr]{1.2.3}).
\end{proof}

\begin{env}[1.4.4]
\label{III.1.4.4-fr}
Si $X$ est un \emph{schéma affine}, il résulte de
(\hyperref[III.1.2.5-fr]{1.2.5}) et
(\hyperref[III.1.4.1-fr]{1.4.1}) que pour tout $q\geq0$, les diagrammes
\begin{equation}
\label{III.1.4.4.1-fr}
\tag{1.4.4.1}
\xymatrix{
H^q(U,\mathcal F'')\ar[r]^{\partial}\ar[d]
&H^{q+1}(U,\mathcal F')\ar[d]\\
H^{q+1}((\mathbf f),M'')\ar[r]^{\partial}
&H^{q+2}((\mathbf f),M')
}
\end{equation}
correspondant à une suite exacte
$0\to\mathcal F'\to\mathcal F\to\mathcal F''\to0$ de
$\mathcal O_X$-Modules quasi-cohérents (avec les notations de
(\hyperref[III.1.2.5-fr]{1.2.5})), sont commutatifs.
\end{env}

